%% =====================================================================
%%  B6-T-23-04 -- what B6 contributes to "The Gut Verdict"
%%  -------------------------------------------------------------------
%%  LAYER A: APPEND-ONLY. Written once at the entry's write, then left
%%  alone. If the source has more to say later, add a bullet -- do not
%%  rewrite. Material found ONLY here goes in the `unique` slot with
%%  @unique set to yes, which makes it undeletable (rule R10).
%% =====================================================================
%% @kind: source
%% @id: B6-T-23-04
%% @book: B6
%% @topic: T-23-04
%% @locator: ch. 9 intro, p. 111; ch. 10 (vulnerabilities), pp. 128--130
%% @status: distilled
%% @unique: yes
%% @integrated: yes
%% @updated: 2026-09-19

\dossier{B6}{T-23-04}{\bookshort{B6}}{ch. 9 intro, p. 111; ch. 10 (vulnerabilities), pp. 128--130}

\begin{distilled}
  \item Anyone dealing with a covertly aggressive person needs heightened gut-level sensitivity to the tactics to avoid being taken in; the tactics are engineered to be deniable at the level of argument while leaving felt discomfort behind.
  \item B6's patients' common report: I never really trusted them but could not pinpoint why --- the alarm is real data, running ahead of analysis.
  \item The override codebook (the target-side vulnerabilities): naivete (refusing to believe anyone is as cunning as the gut reports), over-conscientiousness (pre-crediting their side), low self-confidence (doubting your right to judge), over-intellectualisation (believing understood reasons would excuse behaviour), emotional dependency.
  \item The training order: know all the tactics by heart; watch and listen carefully --- listen for, not necessarily to; and never explain the alarm away before auditing it.
\end{distilled}
\begin{unique}
  \item The vulnerability checklist (naivete, over-conscientiousness, low self-confidence, over-intellectualisation, emotional dependency) and the listen-for-not-to rule are unique to B6 in this corpus.
\end{unique}
